\documentclass[12pt]{article}
%^ Start of document
\usepackage{times}  %Makes text TimesNewRoman
\usepackage{setspace} %Allows you to set single or double space 
\usepackage{biblatex} %Imports biblatex package
\usepackage{graphicx} % Required for inserting images
\usepackage{authblk}
\usepackage[colorlinks=true, urlcolor=blue, linkcolor=blue]{hyperref}
\usepackage{subfig}
\usepackage{float}
\usepackage{tabularx}
\usepackage{multirow}
\usepackage[paperheight=11in,
   paperwidth=8.5in,
   top=20mm,
   bottom=20mm,
   left=40mm,
   right=40mm]{geometry}
\usepackage{moresize}
\usepackage{tikz}
\usepackage{xcolor}
\usepackage{physics}
\usepackage{amssymb}
\usepackage{mathrsfs}
\usepackage{amsmath}
\usepackage[document]{ragged2e}
\usepackage{comment}
\usepackage{xcolor}
\usepackage{amsmath, amssymb}
\usepackage{physics}
\usepackage{bm}
\usepackage{geometry}
\usepackage{mathtools}
\usepackage{tcolorbox}
\addbibresource{Qhw.bib}
\begin{document}
\singlespacing  %Sets single spacing for whole document
\title{Quantum CH 4 HW}

\author[1]{Zachary Tullier}
\affil[1]{Student\\Physics Department\\ University of Houston - Clear Lake \\Houston, Texas\\U.S}
\date{}
\maketitle

\section{Textbook Question 4.11}


\section*{Exercise 4.11}
    Consider the dimensionless Hamiltonian
    \[
        \hat{H} = \frac{1}{2} \hat{P}^2 + \frac{1}{2} \hat{X}^2, \quad
    \]
    with 
    \[
        \hat{P} = -i \frac{d}{dx}
    \]
    
    \begin{enumerate}
        \item[(a)] Show that the wave functions $\psi_0(x) = \frac{e^{-x^2/2}}{\sqrt[4]{\pi}}$ and $\psi_1(x) = \sqrt{\frac{2}{\sqrt{\pi}}} \, x e^{-x^2/2}$ are eigenfunctions of $\hat{H}$ with eigenvalues \( 1/2 \) and \( 3/2 \), respectively.
        \item[(b)] Find the values of the coefficients \( \alpha \) and \( \beta \) such that
        \[
            \psi_2(x) = \frac{1}{\sqrt{2 \sqrt{\pi}}} \left( \alpha x^2 - 1 \right) e^{-x^2/2}
            \quad \text{and} \quad
            \psi_3(x) = \frac{1}{\sqrt{6 \sqrt{\pi}}} \, x \left( 1 + \beta x^2 \right) e^{-x^2/2}
        \]
        are orthogonal to \( \psi_0(x) \) and \( \psi_1(x) \), respectively. Then show that \( \psi_2(x) \) and \( \psi_3(x) \) are eigenfunctions of \( \hat{H} \) with eigenvalues \( 5/2 \) and \( 7/2 \), respectively.
    \end{enumerate}

    \subsection{Solution, part (a):}
        Apply the general harmoic oscillator hamiltonian operator to a wave function
        \[
            \hat{H} \psi_n(x) = E_n \psi_n(x)
        \]
        
        The solutions to this equation are well known and given by
        \[
            \psi_n(x) = \frac{1}{\sqrt{2^n n!}} \sqrt{\frac{1}{\sqrt{\pi}}} H_n(x) e^{\frac{-x^2}{2}} 
        \]

        Where $H_n(x)$ is the Hermite polynomial of degree n and $E_n$ are the energy levels
        The ground state ($n=0$, $H_0(x)=1$) and the first excited state ($n=1$, $H_1(x)=2x$) are
        \[
            \psi_0(x) = \frac{1}{\pi^{\frac{1}{4}}} e^{\frac{-x}{2}}
        \]
        \[
            \psi_1(x) = \frac{1}{\sqrt{2}} \cdot \frac{1}{\pi^{1/4}} \cdot 2x \, e^{-x^2/2} = \sqrt{2} \cdot \frac{x}{\pi^{1/4}} e^{-x^2/2} = \sqrt{\frac{2}{\sqrt{\pi}}} \, x \, e^{-x^2/2}
        \]

        Rescaling the general harmonic oscillator Hamiltonian equation (\ref{eq: hoh}) by simply defining $\hbar = m = \omega = 1$
        \[
            \hat{H} = \frac{1}{2} \left( - \frac{d^2}{dx^2} + x^2 \right)
        \]

        Apply this to both the ground state and the first excited state and simplify
        \[
            \boldsymbol{\boxed{\hat{H} \psi_0 = \frac{1}{2} \psi_0}}
        \]
        \[
            \boldsymbol{\boxed{\hat{H} \psi_1 = \frac{3}{2} \psi_1}}
        \]

    \subsection{Solution, part (b):}        
        We are given 
        \[
            \psi_2(x) = \frac{1}{\sqrt{2 \sqrt{\pi}}} \left( \alpha x^2 - 1 \right) e^{-x^2/2}
        \]
        \[
            \psi_3(x) = \frac{1}{\sqrt{6 \sqrt{\pi}}} \, x \left( 1 + \beta x^2 \right) e^{-x^2/2}
        \]
        Proving that ground and excited state are orthogonal to these given states requires
        \begin{equation} \label{eq: psi0}
            \langle \psi_0 | \psi_2 \rangle = 0
        \end{equation}
        \begin{equation} \label{eq: psi1}
            \langle \psi_1|\psi_3 \rangle = 0
        \end{equation}

        The notations $\langle|\rangle$ are the following general equations
        \[
            \langle \psi_n | \psi_m \rangle = \int_{-\infty}^{\infty} \psi_n(x) \psi_m(x) dx
        \]

        Apply to equation \ref{eq: psi0}
        \begin{equation}
            \begin{split}
                \langle \psi_0 | \psi_2 \rangle = \int_{-\infty}^{\infty} \psi_0(x) \psi_2(x) dx \\
                = \int_{-\infty}^{\infty} \left( \frac{1}{\pi^{1/4}} e^{-x^2/2} \right) \left( \frac{1}{\sqrt{2\sqrt{\pi}}} (\alpha x^2 - 1) e^{-x^2/2} \right) \\
                = \frac{1}{\pi^{1/4} \sqrt{2\sqrt{\pi}}} \int_{-\infty}^{\infty}  (\alpha x^2 - 1) e^{-x^2} dx \\
                = \frac{1}{\sqrt{2} \pi^{3/4}} \left[ \alpha \int_{-\infty}^{\infty} x^2 e^{-x^2} dx - \int_{-\infty}^{\infty} e^{-x^2} dx \right] \\
                = \frac{1}{\sqrt{2} \pi^{3/4}} \left[ \alpha \cdot \frac{\sqrt{\pi}}{2} - \sqrt{\pi} \right] \\
                = \frac{\sqrt{\pi}}{\sqrt{2} \pi^{3/4}} \left( \frac{\alpha}{2} - 1 \right) \\
                = \frac{1}{\sqrt{2} \pi^{3/4}} \left[ \alpha \int_{-\infty}^{\infty} x^2 e^{-x^2} \, dx - \int_{-\infty}^{\infty} e^{-x^2} \, dx \right]
            \end{split}
        \end{equation}

        Use well known gaussian integrals and set equal to 0
        \[
            \int_{-\infty}^{\infty} e^{-x^2}dx = \sqrt{\pi}
        \]
        \[
            \int_{-\infty}^{\infty} x^2 e^{-x^2}dx = \frac{\sqrt{\pi}}{2}
        \]
        \[
            \langle \psi_0 | \psi_2 \rangle = \frac{1}{\sqrt{2} \pi^{3/4}} \left[ \alpha \cdot \frac{\sqrt{\pi}}{2} - \sqrt{\pi} \right] = \frac{\sqrt{\pi}}{\sqrt{2} \pi^{3/4}} \left( \frac{\alpha}{2} - 1 \right) = 0 \Rightarrow \boldsymbol{\boxed{\alpha = 2}}
        \]

        Apply the notations $\langle|\rangle$ general equation to equation \ref{eq: psi1}
        \begin{equation}
            \begin{split}
                \langle \psi_1 | \psi_3 \rangle = \int_{-\infty}^{\infty} \psi_1(x)\, \psi_3(x)\, dx \\
                = \int_{-\infty}^{\infty} \left( \sqrt{\frac{2}{\sqrt{\pi}}} \, x \, e^{-x^2/2} \right) \left( \frac{1}{\sqrt{6\sqrt{\pi}}} \, x(1 + \beta x^2) \, e^{-x^2/2} \right) dx \\
                = \frac{1}{\sqrt{3} \pi^{3/4}} \int_{-\infty}^{\infty} x^2 (1 + \beta x^2) e^{-x^2} dx \\
                = \frac{1}{\sqrt{3} \pi^{3/4}} \left[ \int_{-\infty}^{\infty} x^2 e^{-x^2} \, dx + \beta \int_{-\infty}^{\infty} x^4 e^{-x^2} \, dx \right]
            \end{split}
        \end{equation}

        Use the same well known gaussian integrals and set equal to 0
        \[
            \langle \psi_1 | \psi_3 \rangle = \frac{1}{\sqrt{3} \pi^{3/4}} \left( \frac{\sqrt{\pi}}{2} + \beta \cdot \frac{3\sqrt{\pi}}{4} \right) = \frac{\sqrt{\pi}}{\sqrt{3} \pi^{3/4}} \left( \frac{1}{2} + \frac{3}{4} \beta \right) = 0 \Rightarrow \boldsymbol{\boxed{\beta = -\frac{2}{3}}}
        \]

\section{Textbook Question 4.18}
    A particle of mass \( m \) in an infinite potential well of length \( a \) has the following initial wave function at \( t = 0 \):
    
    \[
        \psi(x,0) = \frac{2}{\sqrt{7a}} \sin\left( \frac{\pi x}{a} \right) 
        + \sqrt{\frac{6}{7a}} \sin\left( \frac{2\pi x}{a} \right) 
        + \frac{2}{\sqrt{7a}} \sin\left( \frac{3\pi x}{a} \right)
    \]
    
    \begin{enumerate}
        \item[(a)] If we measure energy, what values will we find and with what probabilities? Calculate the average energy.
    
        \item[(b)] Find the wave function \( \psi(x,t) \) at any later time \( t \). Determine the probability of finding the particle at a time \( t \) in the state
        \[
            \varphi(x,t) = \frac{1}{\sqrt{a}} \sin\left( \frac{3\pi x}{a} \right) \exp\left( -i \frac{E_3 t}{\hbar} \right).
        \]
    
        \item[(c)] Calculate the probability density \( \rho(x,t) \) and the current density \( \vec{J}(x,t) \). Verify that
        \[
            \frac{\partial \rho}{\partial t} + \nabla \cdot \vec{J}(x,t) = 0.
        \]
    \end{enumerate}

    \subsection{Solution, part (a):}
        Given
        \begin{equation} \label{eq: g}
            \psi(x, 0) = \frac{2}{\sqrt{7a}} \sin\left( \frac{\pi x}{a} \right) +
            \sqrt{\frac{6}{7a}} \sin\left( \frac{2\pi x}{a} \right) +
            \frac{2}{\sqrt{7a}} \sin\left( \frac{3\pi x}{a} \right)
        \end{equation}

        The normalized eigenfunctions of an infinite well are (See appendix section \ref{app A}):
        \[
            \phi_n(x) = \sqrt{\frac{2}{a}} \sin\left( \frac{n \pi x}{a} \right),
            \quad
            E_n = \frac{n^2 \pi^2 \hbar^2}{2ma^2}
        \]
        
        The Wavefunction at $t=0$ is expressed as a linear combination because it is a superposition:
        \[
            \Psi(x,0) = \sum c_n\psi_n(x) = c_1 \psi_1(x) + c_2 \psi_2(x) + c_3 \psi_3(x)
        \]

        Comparing the coefficients of the superposition and the given wave function (equ \ref{eq: g}) by matching forms
        \[
            c_1 \psi_1(x) = c_1 \sqrt{\frac{2}{a}} \sin \left( \frac{\pi x}{a} \right) = \sqrt{\frac{2}{7a}} \sin \left( \frac{\pi x}{a} \right)
        \]
        \[
            c_2 \psi_2(x) = c_2 \sqrt{\frac{2}{a}} \sin\left( \frac{2\pi x}{a} \right) = \sqrt{\frac{6}{7a}} \sin\left( \frac{2\pi x}{a} \right)
        \]
        \[
            c_3 \psi_3(x) = c_3 \sqrt{\frac{2}{a}} \sin\left( \frac{3\pi x}{a} \right) = \sqrt{\frac{2}{7a}} \sin\left( \frac{3\pi x}{a} \right)
        \]

        Simplifying we get
        \[
            c_1 = \sqrt{\frac{2}{7a}} \div \sqrt{\frac{2}{a}} = \sqrt{\frac{4}{14}}
        \]
        \[
            c_2 = \sqrt{\frac{6}{7a}} \div \sqrt{\frac{2}{a}} = \sqrt{\frac{3}{7}}
        \]
        \[
            c_3 = \sqrt{\frac{2}{7a}} \div \sqrt{\frac{2}{a}} = \sqrt{\frac{4}{14}}
        \]

        Using the uncertainty principle:
        If a system is in state ($\Psi$), and you measure an observable, the probability of obtaining the eigenvalue ($E_n$) is
        \begin{equation} \label{eq: P}
            P(E_n) = |c_n|^2       
        \end{equation}
        
        Probabilities:
        \[
            P_1 = \frac{4}{14} = \frac{2}{7}, \quad
            P_2 = \frac{3}{7}, \quad
            P_3 = \frac{4}{14} = \frac{2}{7}
        \]
        
        The average energy equation becomes
        \[
            \boldsymbol{\boxed{\langle E \rangle = \frac{2}{7} E_1 + \frac{3}{7} E_2 + \frac{2}{7} E_3
            = \frac{\pi^2 \hbar^2}{2ma^2} \cdot \frac{2(1^2) + 3(2^2) + 2(3^2)}{7}
            = \frac{32 \pi^2 \hbar^2}{7ma^2}}}
        \]

    \subsection{Solution, part (b):}
        Each eigenstate evolves as:
        \[
            \psi(x,t) = \sum_{n=1}^{3} c_n \phi_n(x) e^{-i E_n t / \hbar}
        \]
        
        \[
            = \frac{2}{\sqrt{14}} \phi_1(x) e^{-i E_1 t / \hbar} +
            \sqrt{\frac{3}{7}} \phi_2(x) e^{-i E_2 t / \hbar} +
            \frac{2}{\sqrt{14}} \phi_3(x) e^{-i E_3 t / \hbar}
        \]
        
        Using probability equation 
        \[
            \varphi(x,t) = \phi_3(x) e^{-i E_3 t / \hbar} \Rightarrow P = \left| \braket{\varphi}{\psi} \right|^2 = \left| \frac{1}{\sqrt{7}} \right|^2 = \frac{1}{7}
        \]
        
    \subsection{Solution, part (c):}
        The probability equation is 
        \[
            \rho(x,t) = |\psi(x,t)|^2 
        \]
        
        \[
            \vec{J}(x,t) = \frac{\hbar}{2 m i} \left( \psi^* \frac{\partial \psi}{\partial x} - \psi \frac{\partial \psi^*}{\partial x}\right)
        \]
        
        This is the standard quantum continuity equation
        \[
            \frac{\partial \rho}{\partial t} + \frac{\partial J}{\partial x} = 0
        \]

        This identity always holds for any wavefunction evolving via the Schrodinger equation.
        

\section{Textbook Question 4.26}
    A particle of mass \( m \) is subject to a delta potential
    \[
        V(x) = 
        \begin{cases}
        \infty, & x \leq 0, \\
        V_0 \delta(x - a), & x > 0,
        \end{cases}
    \]
    where \( V_0 > 0 \).
    
    \begin{enumerate}
        \item[(a)] Find the wave functions corresponding to the cases \( 0 < x < a \) and \( x > a \).
        \item[(b)] Find the transmission coefficient.
    \end{enumerate}

    \subsection{Solution, part (a):}
        A particle of mass \( m \) is subject to the potential:
        \[
            V(x) =
            \begin{cases}
            \infty, & x \leq 0 \\
            V_0 \delta(x - a), & x > 0
            \end{cases}
            \quad \text{where } V_0 > 0
        \]
               
        In the regions where \( V(x) = 0 \), the Schrödinger equation is:
        \[
            \frac{d^2 \psi}{dx^2} + k^2 \psi = 0, \quad \text{with } k = \frac{\sqrt{2mE}}{\hbar}
        \]
        
        General solutions:
        \[
            \psi_1(x) = A \sin(kx), \quad 0 < x < a
        \]
        \[
            \psi_2(x) = B e^{ikx}, \quad x > a
        \]
        
        Apply boundary conditions ($x=a)$:
        \[
        \psi_1(a) = \psi_2(a) \Rightarrow A \sin(ka) = \boldsymbol{\boxed{B e^{ika}}}
        \]    
        
        Derivative discontinuity from delta function:
        \[
            \psi_2'(a) - \psi_1'(a) = \frac{2m V_0}{\hbar^2} \psi(a)
        \]
        \[
            ik B e^{ika} - A k \cos(ka) = \boldsymbol{\boxed{\frac{2m V_0}{\hbar^2} A \sin(ka)}}
        \]
        

    \subsection{Solution, part (b):}
        Assume an incoming wave to the left of the delta and transmitted to the right:
        \[
            \psi(x) =
            \begin{cases}
                Asin(kx) & 0 < x < a \\
                B e^{ikx}, & x > a
            \end{cases}
        \]
        
        Apply boundary conditions at \( x = a \) and continuity:
        \[
            \psi_1(a) = \psi_2(a) \Rightarrow Asin(ka) = Be^{ika}
        \]
        
        Derive discontinuity
        \[
            \frac{\partial \psi_1}{\partial x} = A k cos(kx)
        \]
        \[
            \frac{\partial \psi_2}{\partial x} = i k B e^{ikx}
        \]
        
        Apply jump
        \[
            \psi_{II}'(a) - \psi_{I}'(a) 
            = i k B e^{ika} - A k \cos(ka) 
            = \frac{2m V_0}{\hbar^2} \psi(a) 
            = \frac{2m V_0}{\hbar^2} A \sin(ka)
        \]

        Solve for transmission coefficient ($T$)
        \[
            T = \left| \frac{B}{A} \right|^2 = \frac{|B|^2}{|A|^2} = \frac{A^2 \sin^2(ka)}{|A|^2} = sin^2(ka)
        \]

        Now we substitute back the value of $\frac{B}{A}$ into the jump condition
        \[
            i k B e^{ika} - A k \cos(ka) = \frac{2m V_0}{\hbar^2} A \sin(ka)
        \]
        
        Substitute \( B = A \sin(ka) e^{-ika} \):
        \[
            i k A \sin(ka) - A k \cos(ka) = \frac{2m V_0}{\hbar^2} A \sin(ka)
        \]

        Divide both sides by \( A \sin(ka) \):
        \[
            i k - k \cot(ka) = \frac{2m V_0}{\hbar^2}
        \]      

        Rewriting:
        \[
            \cot(ka) = i - \frac{2m V_0}{\hbar^2 k}
        \]

        And from the continuity condition earlier, we had:
        \[
            \boldsymbol{\boxed{T = \sin^2(ka)}}
        \]

\section{Textbook Question 4.29}
    Consider a particle of mass \( m \) that is moving under the influence of an attractive delta potential
    \[
        V(x) =
        \begin{cases}
        - V_0 \delta(x), & x > -a, \\
        \infty, & x < -a,
        \end{cases}
    \]
    
    where \( V_0 > 0 \). Discuss the existence of bound states in terms of \( V_0 \) and \( a \).

    \subsection{Solution:}
        Consider a particle of mass \( m \) under the potential:
        \[
            V(x) =
            \begin{cases}
            - V_0 \delta(x), & x > -a \\
            \infty, & x < -a
            \end{cases}
            \quad \text{with } V_0 > 0
        \]
        
        We are to discuss the existence of bound states in terms of \( V_0 \) and \( a \).
    
        We assume a bound state \( E = -\frac{\hbar^2 \kappa^2}{2m} \), with \( \kappa > 0 \). In both regions where \( V(x) = 0 \), the Schrödinger equation becomes:
        
        \[
            \frac{d^2 \psi}{dx^2} = \kappa^2 \psi
        \]
        
        Region I: \( -a < x < 0 \)
        \[
            \psi_I(x) = A \sinh[\kappa(x + a)]
            \quad \Rightarrow \psi(-a) = 0
        \]
        
        Region II: \( x > 0 \)
        \[
            \psi_{II}(x) = B e^{-\kappa x}
            \quad \text{(decays at infinity)}
        \]
        
    Boundary Conditions at \( x = 0 \)
        \begin{itemize}
            \item Continuity:
            \[
                \psi_I(0) = \psi_{II}(0) \Rightarrow A \sinh(\kappa a) = B
            \]
        
            \item Derivative discontinuity:
            \[
                \psi'(0^+) - \psi'(0^-) = -\frac{2m V_0}{\hbar^2} \psi(0)
            \]
            \[
                -\kappa B - A \kappa \cosh(\kappa a) = -\frac{2m V_0}{\hbar^2} A \sinh(\kappa a)
            \]
        \end{itemize}
        
        Substitute \( B = A \sinh(\kappa a) \):
        \[
            - \kappa A \left[ \sinh(\kappa a) + \cosh(\kappa a) \right] = -\frac{2m V_0}{\hbar^2} A \sinh(\kappa a)
        \]
        
        \[
            \Rightarrow \kappa \left[ 1 + \coth(\kappa a) \right] = \frac{2m V_0}{\hbar^2}
        \]


\section{Appendix}
    \subsection{Normalized Eigen Functions of an infinite well} \label{app A}
        We are solving the time-independent Schrodinger equation (TISE) in one dimension for a particle of mass $m$ in a potential well of width $a$, with infinite walls
        \[
            V(x) = 
            \begin{cases}
            0 & \text{if } 0 < x < a \\
            \infty & \text{otherwise}
            \end{cases}
        \]
        Therefore at the boundaries
        \[
            \psi(0) = \psi(a) = 0
        \]
        The TISE equation 
        \[
            -\frac{\hbar^2}{2 m} \frac{d^2 \psi(x)}{dx^2} = E \psi(x)
        \]
        Then becomes
        \[
            0 = \frac{d^2 \psi(x)}{dx^2} + k^2 \psi(x)
        \]
        Where 
        \begin{equation} \label{eq: k}
            k = \frac{\sqrt{2 m E}}{\hbar}
        \end{equation}
        The general solution to this second order differential equation is 
        \[
            \psi(x) = A sin(kx) + B cos(kx)
        \]
        We apply the boundary conditions ($x=0$, $x=a$)
        \[
            \psi(0) = A sin(0) + B cos(0) = B = 0
        \]
        \[
            \psi(a) = A sin(ka) \Rightarrow ka=n\pi
        \]
        Where 
        \begin{equation} \label{eq: k2}
            k=\frac{n \pi}{a}    
        \end{equation}
        and $n=1,2,3,...$. So the allowed solutions are 
        \[
            \psi_n(x) = A sin(\frac{n \pi x}{a})
        \]
        Normalize the wave function to get the eigenfunction
        \[
            \int_0^a |\psi_n(x)|^2dx = 1 \Rightarrow A^2 \int_0^a sin^2 \left(\frac{n \pi x}{a} dx = 1 \right)
        \]
        Simplify
        \[
            \int_0^a sin^2 \left( \frac{n \pi x}{a} \right) dx = \frac{a}{2}
        \]
        \[
            A^2 \frac{a}{2} = 1 \Rightarrow A = \sqrt{\frac{2}{a}}
        \]
        The eigenstates of the infinite square well are:
        \[
            \phi_n(x) = \sqrt{\frac{2}{a}} \sin\left( \frac{n\pi x}{a} \right), \quad n=1,2,3,\ldots
        \]
        From equation \ref{eq: k}, we isolate $E$ and substitute equation \ref{eq: k2} getting the eigenstates energies
        \begin{equation} \label{eq: E}
            E = \frac{n^2 \hbar^2 k^2}{2 m a^2}
        \end{equation}


        
\end{document}
